\chapter{Introduction}

The European continent, characterized by its high density of sovereign states and highly integrated, cross-border public transportation infrastructure, presents a unique theater for complex logistical optimization. This thesis tackles a highly constrained, multi-objective optimization problem framed by a distinct real-world application: computationally determining the global optimal routing sequence to break the Guinness World Record for the most countries visited by scheduled public transport within a strict 24-hour window. While superficially resembling a standard shortest-path query, this challenge requires maximizing the accumulation of unique sovereign territories while strictly minimizing temporal transit costs under a hard 24-hour deadline. Because transit connections (combining short-haul aviation and high-speed rail) do not exist continuously but are rigidly bound to specific departure and arrival schedules, traditional static spatial graphs are fundamentally inadequate. Modeling this dimension of time across a densely connected continental network introduces a massive combinatorial explosion, rendering standard brute-force enumeration computationally intractable.

\medskip

To solve this, the research was structured across three major computational phases: data engineering, architectural comparative analysis, and distributed heuristic optimization. The first phase required the construction of a deterministic, time-accurate transportation dataset. Because no unified, free API exists for all European rail and flight networks, a data acquisition pipeline was engineered using a combination of flight REST APIs and programmatic web scraping of centralized rail aggregators. The resulting raw data underwent rigorous normalization to resolve naming inconsistencies, impute missing chronological boundaries, and map all localized schedules to a universally synchronized UTC timeline. To ensure the graph remained traversable, the network was geographically filtered using data-driven frequency heuristics and Large Language Models (LLMs) to isolate only the most critical international transit hubs, while specific intra-city transfers were manually validated to guarantee geospatial accuracy.

\medskip

With the dataset normalized, the second phase focused on evaluating the optimal mathematical environment for the search algorithm. The thesis directly contrasts two fundamental paradigms for modeling time-dependent networks. The initial approach utilized an Activity-on-Edge (AoE) architecture, which maintained a highly compact memory footprint by compressing temporal schedules into dynamic arrays on spatial edges. However, searching this graph required Dynamic Programming (DP) and the constant maintenance of expanding Pareto frontiers. This approach ultimately suffered a computational "time wall", choking on its own comparative loops and proving highly resistant to distributed parallelization due to Inter-Process Communication (IPC) latency. Consequently, the model was redesigned as an Activity-on-Node (AoN) framework. By transforming specific transit events into the graph's vertices, time was explicitly discretized, creating a strictly chronological Directed Acyclic Graph (DAG).

\medskip

This structural pivot unlocked the final phase of the research: large-scale distributed optimization. The chronological nature of the AoN DAG decoupled the search space, allowing a Depth-First Search (DFS) integrated with Branch and Bound (B\&B) pruning to be distributed across a 100-node High-Performance Computing (HPC) cluster. However, initial unconstrained executions revealed a blind spot: the algorithm's pursuit of mathematical minimums led it to exploit physically impossible transfers, which hijacked the B\&B bounding logic and systematically pruned realistic, human-viable routes. To extract the true feasible optimum, a custom heuristic sorting mechanism was engineered to manipulate the algorithm's traversal stack, prioritizing realistic layover times over very short, physically impossible transfers. Finally, by applying a bespoke Route Quality Index (RQI) to score the physical viability of the generated permutations, the pipeline successfully filtered out millions of theoretical but infeasible routes to isolate the exact, physically executable global optimum, definitively establishing the theoretical and human-feasible limits of the European transit network.

The complete, reproducible Python codebase for data acquisition, graph construction, and distributed HPC execution developed for this research is publicly available at: \url{https://github.com/aniol-petit/route-planner-tfg.git}.